% Salin gambar PNG/SVG ke folder image/BabVI sebelum kompilasi.
Analisis GSM8K menggunakan seluruh soal yang tersedia sampai 1.319 soal. Delta terhadap
kondisi netral dihitung secara deskriptif. Delta Man merupakan rata-rata akurasi satu
kondisi eksplisit dan tiga kondisi implisit Man dikurangi akurasi Neutral; Delta Woman
dihitung dengan cara yang sama. Untuk uji inferensial, setiap kondisi persona dibandingkan
dengan Neutral pada nomor soal yang sama menggunakan \textit{two-sided Exact McNemar}.
Nilai $p$ dikoreksi dengan prosedur Holm untuk delapan perbandingan di dalam setiap model.
McNemar tidak diterapkan langsung pada Delta Man atau Delta Woman karena keduanya merupakan
rata-rata beberapa kondisi, bukan satu pasangan keluaran biner.

\begin{figure}[H]
    \centering
    \includegraphics[width=\textwidth]{image/BabVI/01_delta_avg_vs_neutral_all_available.png}
    \caption{Perubahan rata-rata akurasi kondisi persona Man dan Woman terhadap kondisi
    Neutral pada GSM8K. Delta dihitung menggunakan definisi yang sama dengan grafik
    \textit{Avg vs Neutral} pada dokumen analisis awal. Tanda bintang menunjukkan bahwa
    sedikitnya satu perbandingan granular penyusun terhadap Neutral signifikan berdasarkan
    \textit{two-sided Exact McNemar} setelah koreksi Holm ($p_\mathrm{adj}<0{,}05$),
    bukan uji langsung atas delta agregat.}
    \label{fig:gsm8k-delta-avg-vs-neutral}
\end{figure}

Gambar~\ref{fig:gsm8k-delta-avg-vs-neutral} menunjukkan perubahan deskriptif terbesar
pada OLMo 3 32B Think untuk persona Man
($\Delta=+0.051$).
Di antara enam model yang dilanjutkan, 1 model memiliki sedikitnya satu perbandingan yang tetap signifikan setelah koreksi Holm. Hasil ini dibaca bersama
kriteria operasional dan keragaman model, bukan sebagai peringkat mutu model secara umum.
Kimi K2 Thinking dan Trinity Large Thinking dikeluarkan dari grafik karena hanya memiliki
paling banyak 400 soal. Jumlah observasi dan pasangan aktual untuk 14 model yang dianalisis
dilaporkan pada tabel cakupan data dan tabel Exact McNemar.
